% Phần thành viên C: ghép vào báo cáo chung; cần gói amsmath, amssymb.
% !TeX root = main.tex

\section{Giải thuật di truyền và giải thuật lai}
\subsection{Biểu diễn và xử lý ràng buộc}
Mỗi cá thể là véc-tơ nhị phân $x\in\{0,1\}^n$, trong đó $x_j=1$ khi chọn camera $j$. Chi phí và số camera phủ nút $i$ là
\[
 f(x)=\sum_{j=1}^{n}c_jx_j,\qquad q_i(x)=\sum_{j:i\in S_j}x_j.
\]
Lời giải hợp lệ khi $q_i(x)\geq 1$ với mọi $i$. Sau lai ghép và đột biến, xét từng nút chưa phủ; thêm camera phủ nút đó có tỉ số $c_j/|S_j\cap U|$ nhỏ nhất, với $U$ là tập nút chưa phủ hiện tại. Sau đó xét các camera theo chi phí giảm dần, xóa camera nếu mọi nút của nó vẫn được camera khác phủ. Bộ đếm $q_i$ và chi phí được cập nhật gia tăng khi thêm hoặc xóa camera. Với đầu vào khả thi, bước sửa chữa tạo lời giải hợp lệ; xóa thừa vẫn bảo toàn tính hợp lệ.

\subsection{GA}
GA là phương pháp tìm kiếm quần thể dựa trên chọn lọc, lai ghép và đột biến. Thuật toán GA cho set cover của Beasley và Chu là tài liệu tham khảo chuyên biệt; cài đặt này là một biến thể độc lập, không tái hiện toàn bộ toán tử trong bài báo đó.

\noindent\textbf{Mã giả tổng quát:} Khởi tạo quần thể; lặp chọn cha mẹ, lai ghép, đột biến, đánh giá và thay thế; trả về cá thể tốt nhất.

\noindent\textbf{Mã giả cho bài toán:}
\begin{enumerate}
 \item Tạo một cá thể rỗng và các véc-tơ ngẫu nhiên; sửa chữa, xóa thừa để thu quần thể hợp lệ.
 \item Giữ lại cá thể tốt nhất từng gặp. Chọn mỗi cha mẹ bằng tournament gồm ba cá thể lấy mẫu có hoàn lại.
 \item Lai đồng đều: tại mỗi vị trí lấy bit từ một trong hai cha mẹ với xác suất bằng nhau. Đảo mỗi bit độc lập với xác suất $p_m$.
 \item Sửa chữa, xóa thừa, đánh giá chi phí; thêm con vào thế hệ mới. Lặp cho tới hết ngân sách thời gian.
\end{enumerate}
Mặc định kích thước quần thể $P=40$, $p_m=1/n$, xác suất lai đồng đều tại mỗi bit là $1/2$, tournament kích thước $3$, và một cá thể ưu tú được giữ lại. $P$ điều khiển số cá thể; $p_m$ điều khiển mức thay đổi gen. Đây là giá trị khởi đầu, không phải tham số đã được tối ưu. Đánh giá độ nhạy bằng cách thay đổi từng tham số trên tập huấn luyện, giữ cố định các tham số khác, seed và ngân sách; báo cáo kết quả trên tập kiểm thử riêng.

\subsection{GA lai tìm kiếm cục bộ}
Biến thể \texttt{GA\_LS} dùng cùng GA, sau sửa chữa thực hiện tìm kiếm cục bộ trên từng cá thể. Với camera đang chọn $a$, xét các nút chỉ được $a$ phủ. Một camera chưa chọn $b$ có thể thay $a$ nếu phủ tất cả các nút đó và $c_b<c_a$. Chọn camera thay thế rẻ nhất trong các ứng viên, cập nhật $q_i$ gia tăng và tiếp tục với camera tiếp theo. Các bước thay thế chỉ giảm chi phí và giữ tính hợp lệ; cuối lượt xóa cột thừa. Thứ tự camera được xáo trộn theo seed. Đây là một lượt tìm kiếm, không khẳng định đạt cực tiểu cục bộ hoàn toàn.

Gọi $L=\sum_j|S_j|$ là số phần tử khác không. Lưu một cá thể cần $O(n+m)$ bộ nhớ; quần thể cần $O(P(n+m)+L)$. Khởi tạo bộ đếm và xóa thừa tốn $O(n+L)$ ngoài chi phí sắp xếp $O(n\log n)$. Sửa chữa có thể quét các cột phủ từng hàng chưa phủ và danh sách hàng của chúng; cận thô cho một cá thể là $O(mL+n\log n)$, chưa tính tìm kiếm cục bộ. Lân cận hoán đổi trong trường hợp xấu có thể xét $O(n^2m)$ phép kiểm tra phủ; cài đặt dùng tập hợp hàng để kiểm tra quan hệ phủ với thời gian trung bình $O(1)$. Thời gian toàn bộ phụ thuộc số cá thể được xử lý trong ngân sách, không có bảo đảm tối ưu.

\subsection{Thiết kế đánh giá}
So sánh GA và GA lai với cùng máy, dữ liệu và ngân sách; dùng ít nhất 10 seed độc lập trên mỗi thể hiện. Dùng ít nhất 10 thể hiện thuộc ít nhất ba mức kích thước, gồm OR-Library và dữ liệu tự sinh. Kiểm tra tất cả lời giải bằng chương trình kiểm tra của nhóm; so sánh thể hiện nhỏ với tối ưu chính xác và thể hiện chuẩn với tối ưu/BKS có nguồn. Ghi best, trung bình, độ lệch chuẩn, thời gian, tỉ lệ hợp lệ, tỉ lệ đạt tối ưu và đường hội tụ. Với giá trị tham chiếu $f^*>0$, báo cáo
\[
 \operatorname{gap}(f)=\frac{|f-f^*|}{f^*}\times 100\%.
\]
Không suy diễn ưu thế thuật toán trước khi có kết quả thực nghiệm. Ngân sách theo thời gian có thể làm kết quả thay đổi dù cùng seed; lưu seed, cấu hình máy, phiên bản chương trình và tham số từng lần chạy.

\subsection{Tài liệu tham khảo}
\begin{enumerate}
 \item J. E. Beasley và P. C. Chu, ``A Genetic Algorithm for the Set Covering Problem,'' \emph{European Journal of Operational Research}, 94(2), 392--404, 1996.
 \item J. H. Holland, \emph{Adaptation in Natural and Artificial Systems}, University of Michigan Press, 1975.
 \item P. Moscato, ``On Evolution, Search, Optimization, Genetic Algorithms and Martial Arts: Towards Memetic Algorithms,'' Caltech Concurrent Computation Program Report 826, 1989.
\end{enumerate}

\subsection{Kiểm thử ban đầu trên ví dụ của đề}
Trên thể hiện 7 hàng và 6 cột trong đề, chạy mỗi thuật toán với 10 seed từ 1 đến 10, ngân sách 0,1 giây cho mỗi lần, quần thể 40 và xác suất đột biến $1/6$. Cả 20 lời giải đều được bộ kiểm tra chấp nhận và đạt chi phí tối ưu $16$; do đó chi phí tốt nhất và trung bình của mỗi thuật toán là $16$, độ lệch chuẩn và gap đều bằng $0$, tỉ lệ đạt tối ưu là $10/10$. Thời gian đo trung bình gồm khởi động tiến trình lần lượt khoảng $0{,}2574$ giây cho GA và $0{,}2502$ giây cho GA lai. Thử nghiệm dùng Python 3.14.7 trên Windows; tệp kết quả và log hội tụ được lưu trong \texttt{results/member\_c}. Ví dụ này quá nhỏ để kết luận về chất lượng tương đối hoặc độ nhạy tham số; cần thực nghiệm trên các thể hiện và ngân sách do nhóm thống nhất.

\subsection{Kiểm thử chéo trên dữ liệu tự sinh nhỏ}
Sau khi thành viên B bổ sung dữ liệu tự sinh, chạy kiểm thử chéo trên bốn thể hiện nhỏ với 10 seed từ 1 đến 10, ngân sách 0,05 giây, $P=40$ và $p_m=1/n$. Cả 80 lời giải của GA và \texttt{GA\_LS} đều được bộ kiểm tra độc lập chấp nhận. Kết quả chi phí như sau:
\begin{center}
\begin{tabular}{lrrrrr}
\hline
Thể hiện & Đối chiếu & GA best & GA avg & GA\_LS best & GA\_LS avg \\
\hline
\texttt{nho\_10x20\_s101} & 65  & 65  & 65  & 65  & 65 \\
\texttt{nho\_15x25\_s102} & 136 & 136 & 136 & 136 & 136 \\
\texttt{nho\_20x28\_s103} & 212 & 212 & 212 & 212 & 212 \\
\texttt{nho\_25x30\_s104} & 133 & 133 & 133 & 133 & 133 \\
\hline
\end{tabular}
\end{center}
Độ lệch chuẩn bằng 0 ở cả tám tổ hợp thể hiện--thuật toán. Các giá trị đối chiếu do thành viên B cung cấp; chỉ xem là tối ưu chính thức sau khi thuật toán chính xác của thành viên A xác nhận. Đây là kiểm thử chức năng với ngân sách rất ngắn, không dùng để kết luận GA và GA lai có chất lượng tương đương trên dữ liệu lớn. Script \texttt{experiments/member\_c/run\_tusinh.py} tái lập phép kiểm thử và ghi CSV theo schema kết quả thô của nhóm.

\subsection{Phân tích độ nhạy tham số}

Tham số chính của \texttt{GA} và \texttt{GA\_LS} là kích thước quần thể $P$ và xác suất đột biến $p_m$. Thiết kế thí nghiệm quét một tham số mỗi lần, giữ tham số kia ở giá trị mặc định ($P=40$, $p_m=1/n$), với ngân sách $1{,}0$ giây/lần và 3 seed ($1,2,3$) trên bảy thể hiện: \texttt{nho\_10x20\_s101} ($n=20$), \texttt{vua\_100x400\_s201} ($n=400$), \texttt{scp41} ($n=1000$), \texttt{vua\_150x600\_s202} ($n=600$), \texttt{vua\_200x1000\_s203} ($n=1000$), \texttt{scp51} ($n=2000$), \texttt{scp61} ($n=1000$). Tất cả $420$ lời giải đều được bộ kiểm tra chấp nhận (tỉ lệ hợp lệ $100\%$).

\subsubsection{Độ nhạy theo kích thước quần thể}

$P$ có độ nhạy rõ rệt trên các thể hiện vừa và lớn, nhưng không đáng kể trên thể hiện nhỏ quá dễ.

\begin{center}
\begin{tabular}{lrrrr}
\hline
Thể hiện ($n$) & $P$ & GA best & GA avg $\pm$ std & GA\_LS best \\
\hline
\texttt{nho\_10x20\_s101} (20) & 10 & 65 & $65{,}00 \pm 0{,}00$ & 65 \\
 & 160 & 65 & $65{,}00 \pm 0{,}00$ & 65 \\
\texttt{vua\_100x400\_s201} (400) & 10 & 208 & $218{,}67 \pm 7{,}72$ & 204 \\
 & 160 & 206 & $206{,}00 \pm 0{,}00$ & 204 \\
\texttt{vua\_150x600\_s202} (600) & 10 & 277 & $287{,}33 \pm 7{,}32$ & 261 \\
 & 160 & 261 & $263{,}67 \pm 3{,}77$ & 261 \\
\texttt{vua\_200x1000\_s203} (1000) & 10 & 328 & $336{,}33 \pm 6{,}55$ & 317 \\
 & 80 & 317 & $318{,}00 \pm 1{,}41$ & 317 \\
\texttt{scp41} (1000) & 10 & 432 & $439{,}33 \pm 6{,}60$ & 432 \\
 & 40 & 435 & $435{,}00 \pm 0{,}00$ & 432 \\
 & 160 & 435 & $452{,}00 \pm 12{,}03$ & 434 \\
\texttt{scp51} (2000) & 10 & 267 & $269{,}33 \pm 2{,}05$ & 253 \\
 & 80 & 256 & $258{,}67 \pm 3{,}09$ & 253 \\
 & 160 & 274 & $276{,}33 \pm 1{,}70$ & 268 \\
\texttt{scp61} (1000) & 10 & 144 & $148{,}33 \pm 4{,}78$ & 140 \\
 & 80 & 141 & $141{,}33 \pm 0{,}47$ & 140 \\
\hline
\end{tabular}
\end{center}

Qua bảng trên:
\begin{itemize}
 \item Trên \texttt{vua\_100x400\_s201}, tăng $P$ từ 10 lên 160 làm GA giảm trung bình từ $218{,}67$ xuống $206{,}00$; GA\_LS giảm từ $211{,}33$ xuống $204{,}67$.
 \item Trên \texttt{vua\_150x600\_s202}, $P=10$ cho GA trung bình $287{,}33$ trong khi $P=160$ cho $263{,}67$; GA\_LS ổn định ở $261{,}00$ từ $P=160$.
 \item Trên \texttt{scp51}, $P=80$ tốt hơn cả $P=10$ ($258{,}67$ vs $269{,}33$ cho GA) nhưng $P=160$ lại kém hơn ($276{,}33$), cho thấy dân số quá lớn không phải lúc nào cũng tốt do thời gian cố định bị chia nhỏ cho từng cá thể.
 \item Trên \texttt{scp61} và \texttt{nho\_10x20\_s101}, các mức $P$ đều cho kết quả gần như giống nhau (bài quá nhỏ hoặc quá dễ).
\end{itemize}

\subsubsection{Độ nhạy theo xác suất đột biến}

$p_m$ có độ nhạy thấp. Không có mức nào rõ rệt vượt trội trên mọi thể hiện. Ví dụ trên \texttt{vua\_100x400\_s201}, GA tại $p_m=10\times1/n$ đạt best $204$ (tốt nhất) còn tại $p_m=1/n$ chỉ đạt $211$; nhưng trên \texttt{scp51}, các mức $p_m$ khác nhau chỉ dao động trong khoảng $256$--$264$. Do đó $p_m=1/n$ là lựa chọn mặc định an toàn, không cần chỉnh tinh.

\subsubsection{So sánh GA và GA\_LS}

\texttt{GA\_LS} luôn tốt hơn hoặc bằng \texttt{GA} trên mọi cấu hình tham số và mọi thể hiện. Khoảng cách rõ rệt nhất trên \texttt{vua\_150x600\_s202} ($P=40$: GA avg $276{,}00$ vs GA\_LS avg $268{,}67$) và \texttt{scp51} ($P=80$: GA avg $258{,}67$ vs GA\_LS avg $254{,}33$). Điều này chứng minh việc kết hợp GA với tìm kiếm cục bộ (thuật toán lai) giúp cải thiện chất lượng nghiệm và ổn định hơn so với GA thuần túy.

\subsubsection{Khuyến nghị tham số mặc định}

Dựa trên kết quả trên, tham số mặc định cho \texttt{GA} và \texttt{GA\_LS} được chọn là:
\begin{itemize}
 \item Kích thước quần thể $P = 40$ (điểm cân bằng giữa chất lượng và thời gian trên các thể hiện vừa và lớn).
 \item Xác suất đột biến $p_m = 1/n$ (ổn định, không nhạy).
 \item Kích thước tournament $= 3$, lai ghép đồng đều từng bit với xác suất $1/2$, giữ một cá thể ưu tú (elitism).
 \item Thuật toán khuyến nghị dùng là \texttt{GA\_LS} (tốt hơn \texttt{GA} trên mọi bài kiểm tra).
\end{itemize}

Tất cả script và CSV chi tiết được lưu trong \texttt{experiments/member\_c/run\_ga\_sensitivity.py} và \texttt{results/member\_c/ga\_sensitivity.csv}.
